\subsection{扭余切丛}\label{shifted}

迄今为止，通过考虑 $Y=T^*W$ 及其标准全纯辛结构的情形，我们已能把膜量子化与几何量子化作比较，但只在相当特殊的情形。本节把构造推广到有时称为\textbf{扭余切丛}的对象。\footnote{关于此主题在更一般背景下的更多细节，见 \cite{Ginzburg}。}我们将看到，推广到扭余切丛使膜量子化与几何量子化之间更有趣的比较成为可能。

余切丛 $Y=T^*W$ 有全纯投影 $\pi:T^*W\to W$。纤维是拉格朗日子流形，都是 $\C^n$ 的拷贝，$n=\frac12\dim_\C Y$。此外，$W$ 可作为零截面全纯嵌入 $T^*W$。这个嵌入 $\iota:W\to T^*W$ 是 $\pi:T^*W\to W$ 的“截面”，即 $\pi\circ\iota=1$（等价地，对任意 $w\in W$，$\iota(w)\in\pi^{-1}(w)$）。

下面作推广：考虑具有上述所有性质、只差最后一条的复辛流形 $Y$。换言之，假定有全纯投影 $\pi:Y\to W$，$W$ 为某复流形。纤维是拉格朗日子流形，且同构于 $\C^n$。但投影 $\pi:Y\to W$ 可能没有全纯截面。

我们把这样的结构称为 $Y$ 的\textbf{全纯极化}，一般记作 $\varPi$。$\pi:Y\to W$ 的纤维就是 $\varPi$ 的叶。

对任意全纯线丛 $\L\to W$，我们将构造这样的结构。在 $\L$ 上取酉联络，其曲率 $\sff$ 为 $(1,1)$ 型；局部坐标下 $\sff=-\i\sum_{ij}f_{i\bar j}\d q^i\d\bar q^{\bar j}$。于是 $\Omega_0=\sum_i\d p_i\wedge\d q^i+\sff$ 是 $Y_0=T^*W$ 上的闭 2-形式，当然不是 $(2,0)$ 型。我们将说明如何修改 $Y_0$ 上的复结构，使 $\Omega$ 成为 $(2,0)$ 型且全纯。

描述 $\L$ 的一个显式办法是用开集 $\U_\alpha$ 覆盖 $W$，并在每个开集上选取 $\L$ 的平凡化。在此平凡化下，$\L$ 在 $\U_\alpha$ 上的截面就是全纯函数 $s_\alpha$。在交叠 $\U_\alpha\cap\U_\beta$ 上，局部平凡化由全纯转移函数 $e^{\v_{\alpha\beta}}$ 粘合；$\L$ 的整体截面 $s$ 是一族局部截面 $s_\alpha$，满足
\be\label{comsect} s_\alpha=e^{\v_{\alpha\beta}}s_\beta\ee
于 $\U_\alpha\cap\U_\beta$。$\L$ 上的 Hermite 内积于是在 $\U_\alpha$ 中由 $||s_\alpha||^2=e^{h_\alpha}\bar s_\alpha s_\alpha$ 定义，$h_\alpha$ 为实值函数。要求当 (\ref{comsect}) 成立时 $||s_\alpha||^2=||s_\beta||^2$ 于 $\U_\alpha\cap\U_\beta$，就给出 $h_\alpha$ 与 $h_\beta$ 的关系：
\be\label{hermrel} h_\alpha=h_\beta+2\Re\,\v_{\alpha\beta}. \ee
$\L$ 的曲率在区域 $\U_\alpha$ 中定义为
\be\label{curvdef} \sff=-\i\partial\bar\partial h_\alpha.\ee
式 (\ref{hermrel}) 说明它在交叠区域与 $\alpha$ 无关，故 $\sff$ 整体定义。此外，对 (\ref{hermrel}) 求导得
\be\label{ermrel}\partial_i h_\alpha=\partial_i h_\beta+\partial_i\v_{\alpha\beta}. \ee

我们要作的构造不仅对线丛有效，对线丛的复幂也有效，故先作推广。设 $\L_s$，$s=1,\cdots,k$，是 $W$ 上的全纯线丛，全纯转移函数为 $e^{\v_{s\,\alpha\beta}}$，Hermite 度量由局部函数 $h_{s\,\alpha}$ 定义。为把上一段的陈述推广到 $\L=\otimes_{s=1}^k\L_s^{\c_s}$（$\c_s$ 为复系数），定义曲率形式 $\sff=\sum_{s=1}^k\c_s\sff_s$，并令 $h_\alpha=\sum_{s=1}^k\c_s h_{s\,\alpha}$、$\v_{\alpha\beta}=\sum_s\c_s\v_{s\,\alpha\beta}$。上一段的公式无论 $\c_s$ 是否为整数都成立。特别地，在每个 $\U_\alpha$ 中，(\ref{curvdef}) 与 (\ref{ermrel}) 成立。若 $\c_s$ 不是整数，唯一出问题的地方是：在三重交叠中，$e^{\v_{\alpha\beta}+\v_{\beta\gamma}+\v_{\gamma\alpha}}$ 是常数但未必为 1。这不影响下文。

闭 2-形式 $\Omega=\sum_i\d p_i\d q^i+\sff$ 在 $Y_0=T^*W$ 的自然复结构下当然不是全纯的，甚至不是 $(2,0)$ 型。稍后我们将修改 $Y_0$ 的复结构使 $\Omega$ 全纯。把这个带修改后复结构的流形称为\textbf{扭余切丛}，记作 $Y$。$Y$ 仍带有全纯投影 $\pi:Y\to W$，纤维为 $\C^n$。

在解释此构造前，注意 $\sum_i\d p_i\d q^i=\d(\sum_i p_i\d q^i)$ 是恰当的，故 $\Omega=\sum_i\d p_i\d q^i+\sff$ 的上同调类就是 $\sff$ 的上同调类。这可表述为第一 Chern 类的陈述：若 $[\sff]$ 是形式 $\sff$ 的上同调类，则 $[\sff]/2\pi$ 表示第一 Chern 类 $c_1(\L)$。故
\be\label{jnko} [\Omega]=2\pi c_1(\L)=2\pi\sum_s \c_s c_1(\L_s). \ee
即使全纯投影 $\pi:Y\to W$ 没有全纯截面，它也总有光滑截面。而且 $Y$ 可缩到此截面上，而该截面是 $W$ 的拷贝。故可把式 (\ref{jnko}) 看作 $W$ 的上同调中的陈述。

例如，若 $W$ 紧，则由 Hodge 理论 $c_1(\L)$ 为 $(1,1)$ 型，故式 (\ref{jnko}) 表明 $[\Omega]$ 为 $(1,1)$ 型。又 $c_1(\L_s)$ 对一切 $s$ 为实，故若系数 $\c_s$ 为实，则 $[\Omega]$ 为实。这些是构造有趣量子化例子的便利条件。对量子化，我们要考虑支撑在纤维化 $\pi:Y\to W$ 的非全纯截面 $M$ 上的膜，并要求 $M$ 对 $\Im\Omega$ 为拉格朗日、对 $\Re\Omega$ 为辛。为使这样的 $M$ 存在，条件 $[\Im\Omega]=0$ 是必要的。$[\Re\Omega]$ 非零并非构造例子所必需，但许多最有意思的例子都有 $[\Re\Omega]\ne0$。

为定义新复结构，只须说：在 $T^*\U_\alpha$ 中，函数
\be\label{newc} p_{\alpha\,i}=p_i+\i\partial_i h_\alpha(q,\bar q) \ee
是纤维上的全纯坐标；我们定义新的全纯结构，使 $p_{\alpha\,i}$ 与 $q^j$ 全纯，而非 $p_i$ 与 $q^j$。这个定义已安排好，使得在 $T^*\U_\alpha$ 中
\be\label{ewc} \Omega=\sum_i\d p_{\alpha\,i}\d q^i, \ee
故 $\Omega$ 在 $T^*\U_\alpha$ 中为 $(2,0)$ 型且全纯。当然，$\Omega$ 整体定义，此式对一切 $\alpha$ 成立。此外，由式 (\ref{ermrel})，在交叠区域 $\U_\alpha\cap\U_\beta$ 中
\be\label{newco} p_{\alpha\,i}=p_{\beta\,i}+\i\partial_i\v_{\alpha\beta},\ee
其中 $\v_{\alpha\beta}$ 全纯。于是，虽然 $p_{\alpha\,i}\ne p_{\beta\,i}$，但在 $T^*(\U_\alpha\cap\U_\beta)$ 上使 $p_{\alpha\,i}$ 与 $q^j$ 全纯的复结构，与在该开集上使 $p_{\beta\,i}$ 与 $q^j$ 全纯的复结构相同。因此，我们在各开集 $T^*\U_\alpha$ 上定义的复结构在所有交叠上一致，拼成 $T^*W$ 上的新复结构。这样就定义了新复流形 $Y$，它作为光滑流形与原余切丛相同。在 $Y$ 上，$\Omega$ 全纯且为 $(2,0)$ 型，故 $Y$ 是复辛流形。

式 (\ref{newco}) 与式 (\ref{transpo}) 形式相同，只需取 $w=v/\hbar$，故先前的讨论适用。于是，若能假定量子算符 $\h p,\h q$ 的变换律与经典函数 $p,q$ 的变换律相同，则由第 \ref{cotangent} 节的讨论可推出：$\Hom(\B_\cc,\B_\cc)$ 由作用在 $\L=\otimes_{s=1}^k\L_s^{\lambda_s}$（$\lambda_s=\c_s/\hbar$）上的全纯微分算符组成。$Y$ 的几何由参数 $\c_s$ 控制；若固定 $Y$ 的几何而改变 $\hbar$，则指数 $\lambda_s$ 为 $1/\hbar$ 阶。

这不是正确答案，因为它忽略了第 \ref{cotangent} 节与第 \ref{lagbr} 节分析过的量子效应。即使经典公式 (\ref{newco}) 中 $\v_{\alpha\beta}=0$（例如 $Y=T^*W$ 带其通常辛结构时），也有量子贡献，使得 $\Hom(\B_\cc,\B_\cc)$ 由作用在 $K_W^{1/2}$ 的截面（而非函数）上的全纯微分算符组成。

自然猜测：$\h p$ 变换律中的经典效应与量子效应应当相加。此时 $\Hom(\B_\cc,\B_\cc)$ 将由作用在张量积 $K_W^{1/2}\otimes\L$ 上的全纯微分算符组成。

为论证此答案，只须注意问题有伸缩律：$q$ 不变，而 $p$、$\hbar$、$v$ 均以次数 1 伸缩。故在保持伸缩对称的量子化中，$\h p$ 的变换只能是 $\hbar$ 与 $\phi$ 的线性函数，从而把经典效应与量子效应相加、取 $K_W^{1/2}$ 与 $\L$ 的张量积是正确的。

最后讨论对给定 $W$，如何分类这样构造的 $Y$。设 $W$ 紧。则由 Hodge 理论，任何满足 $c_1(\RR)$ 以复系数为零的全纯线丛 $\RR\to W$ 都容许平坦联络。对此类 $\RR$，作用在 $\L$ 上的全纯微分算符与作用在 $\RR\otimes\L$ 上的相同。因此，若 $W$ 紧，用此构造能作出的复辛流形 $Y$ 完全由复系数的 $c_1(\L)$ 分类，换言之（由式 (\ref{jnko})）由上同调类 $[\Omega]$ 分类。对非紧 $W$，一般不是这样；可以有拓扑平凡但不容许平坦联络的全纯线丛 $\RR\to W$。故非紧情形一般更复杂。
